%% ============================================================
%%  Section 1 — Introduction
%%  Article: BRKGA with Rolling Horizon for the MDRP
%%  Citation style: LNCS (numeric \cite{})
%% ============================================================

\section{Introduction}\label{sec:intro}

The rapid growth of on-demand food delivery platforms has fundamentally
reshaped urban logistics, transforming the way consumers access restaurant
meals and creating a new class of operational challenges for platform
operators~\cite{savelsbergh2022challenges}.
Services such as Uber Eats\textsuperscript{\textregistered},
DoorDash\textsuperscript{\textregistered},
Rappi\textsuperscript{\textregistered},
Just Eat\textsuperscript{\textregistered}, and
Glovo\textsuperscript{\textregistered}
act as digital intermediaries that simultaneously coordinate
customers, restaurants, and crowdsourced couriers across dense
urban environments~\cite{reyes2018mdrp}.
The core operational challenge—assigning couriers to orders and
computing delivery routes in real time—is formally known as the
\emph{Meal Delivery Routing Problem} (MDRP)~\cite{reyes2018mdrp}.

The MDRP belongs to the family of Dynamic Vehicle Routing Problems
(DVRPs) and inherits its NP-hardness from the classical
Vehicle Routing Problem (VRP)~\cite{wen2010dynamic}.
What distinguishes the MDRP from standard VRP variants is the
combination of: (i)~orders that arrive continuously and unpredictably
throughout the day; (ii)~tight customer time windows driven by
freshness expectations; (iii)~heterogeneous courier fleets
with vehicle-dependent speeds; and (iv)~the need to make
assignment and routing decisions within seconds to maintain
operational viability~\cite{ulmer2021restaurant,reyes2018mdrp}.
These characteristics make it impractical to solve the MDRP to
optimality in real time, motivating the development of effective
heuristic and metaheuristic approaches.

Two broad families of solution methods have been explored in
the literature.
\emph{Greedy rolling-horizon heuristics} decompose the planning
horizon into short time windows and apply fast constructive
rules within each window, enabling near-real-time
decision-making at the cost of solution quality~\cite{auad2023courier}.
\emph{Population-based metaheuristics}, such as genetic algorithms
and evolutionary strategies, search a richer solution space and
typically yield higher-quality solutions, but require substantially
more computation time, which can be a limiting factor in dynamic
settings~\cite{liao2020multiobjective}.
Bridging these two paradigms—high solution quality within a
computationally tractable rolling-horizon framework—remains
an open research challenge.

This paper addresses that challenge by proposing a
\emph{Biased Random-Key Genetic Algorithm with Rolling Horizon}
(BRKGA-RH) for the MDRP.
The BRKGA~\cite{goncalves2011brkga} encodes candidate solutions
as vectors of random keys in $[0,1]$ and evolves a population
through biased crossover, so that elite solutions are
preferentially inherited by offspring, providing a principled
balance between intensification and diversification.
The rolling-horizon wrapper re-invokes the BRKGA at every
planning window, allowing the algorithm to react to new orders
and updated courier availability without global
re-optimisation~\cite{reyes2018mdrp}.
We instantiate the framework with two competing objective
functions: (i)~\textbf{BRKGA-Coverage}, which prioritises
maximising the number of orders fulfilled, and
(ii)~\textbf{BRKGA-Travel}, which minimises average travel time
per order while maintaining an acceptable coverage level.

To contextualise the performance of both variants we compare them
against \textbf{MaxCover}, a deterministic greedy rolling-horizon
heuristic that generates all feasible courier--order candidate
assignments within each time window and selects a non-conflicting
subset that maximises coverage.
MaxCover is designed to be a strong fast baseline: it produces
solutions in a fraction of the time required by the BRKGA but
applies no population-level search.
The comparison between the two paradigms on the same set of
instances allows us to quantify the quality--speed trade-off
and to identify the instance characteristics under which
metaheuristic search provides the greatest benefit.

Experiments are conducted on 24 real-world instances derived from
a Latin American food delivery application, spanning small
($<$5{,}000 orders), medium (5{,}000--20{,}000 orders), and
large ($>$20{,}000 orders) operational scenarios.
The dataset includes geo-referenced locations of customers,
restaurants, and couriers, as well as courier working-hour
schedules and vehicle type profiles.

The main contributions of this work are:

\begin{enumerate}
    \item A BRKGA metaheuristic integrated with a rolling-horizon
    framework for the MDRP, with a chromosome encoding that
    guarantees solution feasibility without repair operators.

    \item Two fitness function variants that target, respectively,
    order coverage and travel-time efficiency, enabling
    practitioners to select the objective that best aligns with
    their platform's priorities.

    \item An extensive computational study on 24 heterogeneous
    real-world instances comparing BRKGA-Coverage, BRKGA-Travel,
    and MaxCover across solution quality, travel efficiency,
    and computational cost metrics.

    \item Empirical evidence that the BRKGA variants outperform
    the Business Model Lookahead on small and medium instances
    by an average of 11--13 percentage points in coverage rate,
    while the Business Model Lookahead remains competitive on
    large instances where the bounded generation budget limits
    the metaheuristic's search.
\end{enumerate}

The remainder of the paper is organised as follows.
Section~\ref{sec:litreview} surveys the relevant literature on
the MDRP, rolling-horizon methods, and metaheuristics for
dynamic routing.
Section~\ref{sec:problem} formally defines the problem.
Section~\ref{sec:methodology} describes the MaxCover heuristic
and the proposed BRKGA-RH framework in detail.
Section~\ref{sec:experiments} presents the experimental setup
and the computational results.
Finally, Section~\ref{sec:conclusions} summarises the findings
and outlines directions for future research.
